
To determine how well a geometric model reproduced the polarization observations, a chi-squared value was calculated. First, the difference between the observed and model-predicted polarization angle needed to be calculated carefully. Polarization angles can range from 0--180\degg because polarization measures orientation, rather than direction, meaning that angles that are 180\degg apart have the same polarization vector orientation. As a result, two polarization angles that are physically very similar can have a large numerical difference. For example, polarization angles of 1\degg and 179\degg only differ by 2\deggNS, but directly taking their difference would give 178\deggNS. To account for this, the angular difference is defined by the minimum angular separation between the two angles:
\begin{equation}
    \Delta \theta = \text{min}(|\theta_{\text{obs}} - \theta_{\text{model}}|, 180\deggEQ - |\theta_{\text{obs}} - \theta_{\text{model}}|),
    \label{eq: minimum_separation}
\end{equation}

\noindent where $\theta_{\text{obs}}$ is the observed polarization angle and $\theta_{\text{model}}$ is the polarization angle predicted by either the ratio or Gaussian model. The angular differences were then used to calculate the chi-squared value:
\begin{equation}
    \chi^2 = \sum \left(\frac{\Delta \theta_i}{\sigma_\theta}\right)^2,
    \label{eq: chi_sq}
\end{equation}
\noindent where the angular distance is summed over all sampled polarization vectors. 

The value of $\chi^2$ was then converted to reduced chi-squared by:
\begin{equation}
    \chi^2_{reduced} = \frac{1}{n-k} \times \chi^2,
    \label{eq: chi_sq_reduced}
\end{equation}
\noindent where $n$ is the number of sampled polarization vectors and $k$ is the degrees of freedom. For the ratio model, there is one degree of freedom, and for the Gaussian method, there are three ($\phi$, $b$, and $a$). 